\begin{lemma}\label{lem:decic48cusp13}
Let $X=X_1(48)_{\mathbf F_5}$, and put
$A_{13}=T_{13}-1-13\langle13\rangle$.
For every $\mathbf F_5$-rational cuspidal divisor $C$ on $X$, the divisor
$A_{13}C$ is principal.
\end{lemma}
\begin{proof}
We use primitive cusp labels $(a,b)\in(\mathbf Z/48\mathbf Z)^2$, modulo
$(a,b)\sim\pm(a+nb,b)$. The arithmetic and diamond actions are
\[
 \sigma_u(a,b)=(ua,b),\qquad
 \langle u\rangle(a,b)=(ua,u^{-1}b).
\]
There are $56$ geometric cusps. Their $\sigma_5$-orbits consist of
$16$ rational cusps, $12$ quadratic cusps, and $4$ quartic cusps.
On cusps one has
\[
 T_{13}=\sigma_{13}+13\langle13\rangle\sigma_{13}^{-1}.
\]
Applying this identity to the $32$ closed cusps shows that $A_{13}$ vanishes
on all but four of them. These four form two pairs: the two quadratic
cusps of width $2$, and the two quartic cusps of width $1$.
If their differences are denoted by $C_2-C_2'$ and $C_4-C_4'$, respectively,
the nonzero columns of $A_{13}$ are
\[
 \pm12(C_2-C_2'),\qquad \pm12(C_4-C_4').
\]
This calculation is a permutation calculation on primitive labels;
\path{independent_cusp_A13_labels.json} records every column.

Both displayed differences, multiplied by $12$, are principal. More
precisely, the certificate \path{cusp_corrections_v2.m} gives the two
prime pairs and rational functions $f_2,f_4$ in the degree-$19$
function-field presentation of $X$, with
\[
 \operatorname{div}(f_i)=-12(C_i-C_i')\qquad(i=2,4).
\]
The primes are characterized by their degrees and their pole orders
$2$ and $1$ for $j$, respectively. The identities were checked in both
the finite and infinite maximal orders of a fresh copy of the function
field; see \path{independent_cusp_principal.m}. This proves the claim.
\end{proof}

\begin{lemma}\label{lem:decic48NC13}
Let $D$ be an effective rational divisor of degree $10$ on $X_1(48)$,
and write its reduction at $5$ as $\overline D=N+C$, with $C$ cuspidal
and $N$ supported away from the cusps. Then $A_{13}N$ is principal.
If $C=0$, the same conclusion holds with $13$ replaced by $7$.
\end{lemma}
\begin{proof}
The group $J_1(48)(\mathbf Q)$ is finite by DEHMZ, Theorem~3.1(2).
At either good odd prime $\ell=7,13$, the Eichler--Shimura relation and
injectivity of reduction on rational torsion give
\[
 (T_\ell-1-\ell\langle\ell\rangle)J_1(48)(\mathbf Q)=0.
\]
The same operator annihilates the degree part after subtracting a rational
cusp. Such a cusp is annihilated as a divisor. Reducing at $5$ therefore
gives $A_{13}(N+C)\sim0$. Lemma~\ref{lem:decic48cusp13} gives
$A_{13}C\sim0$, so $A_{13}N\sim0$. If $C=0$, the argument applies
directly to $A_7$.
\end{proof}


\begin{proposition}\label{prop:decic48}
There is no elliptic curve over a number field of degree $10$ with a point
of order $48$.
\end{proposition}
\begin{proof}
Put $X=X_1(48)$ and $Y=X_1(2,24)$. The curve $X$ has genus $37$,
$J_X(\mathbf Q)$ has rank $0$ by \cite[Theorem~3.1(2)]{DEHMZ}, and
$\operatorname{gon}_{\mathbf Q}(X)=16$ by Varivoda, Theorem~D.
Suppose that $D$ is the noncuspidal rational trace divisor of degree $10$
supplied by a point of order $48$.

Its reduction at $5$ is not entirely cuspidal. Otherwise, lift that
reduction to an effective cuspidal divisor $C$ over $\mathbf Q_5$.
The class $[D-C]$ is torsion, by rank zero and Manin--Drinfeld, and has
zero reduction. Torsion injectivity over $\mathbf Q_5$ gives $D\sim C$.
Since the two effective divisors are distinct, base change would give
$h^0(X,D)\geq2$ over $\mathbf Q$, contrary to the gonality bound.
Write $\overline D=N+C$, with $N\ne0$ noncuspidal and $C$ cuspidal.

We first use the degree-$2$ map
\[
 \pi:X\longrightarrow Y,\qquad
 (E,P)\longmapsto
 (E/\langle24P\rangle,P_2,\overline P),
\]
where $P_2$ is the nonzero dual-kernel point. The point $\overline P$
has order $24$ and $P_2\ne12\overline P$. The map preserves cusps.
The Jacobian $J_Y$ has rank $0$, since it is a quotient of $J_X$.
Thus \cite[Theorem~3.1, Proposition~4.16 and Table~3]{DEHMZ} gives
\[
 J_Y(\mathbf Q)=C_Y(\mathbf Q)
 \simeq C_2\times C_4^2\times C_{120}\times C_{240}.
\]
The characteristic-zero cusp blocks are identified by the integral
modular units $b,d,c^2-4d,x(kQ)$ in the normalization
$y^2=x^3+cx^2+dx$, $P_2=(0,0)$, $Q=(b,b)$,
$d=b(1-b-c)$. There are $19$ such blocks, of degrees $1^8,2^6,4^5$;
their integral identification is retained in
\path{inputs/integral_unit_blocks.tex}.

The proven finite Picard presentation gives
\[
 J_Y(\mathbf F_5)\simeq C_2^2\times C_4^2\times C_{12}
                 \times C_{24}\times C_{480}^2.
\]
The rational cusp blocks generate a subgroup $H$ of order $921600$.
Thus $H$ is the full specialization image of $J_Y(\mathbf Q)$,
and $J_Y(\mathbf F_5)/H$ has invariants $[2,2,24,48]$.
Choose the recorded rational cusp $A$ of $Y$, and let $q$ be the
quotient map. If $r=\deg N$, we necessarily have
\[
 -q[\pi_*N-rA]\in\mathcal C_{10-r},\qquad
 \mathcal C_e=\{q[C'-eA]:C'\text{ effective cuspidal of degree }e\}.
 \tag{48.1}
\]
Every finite target cusp is allowed. The cardinalities of these sets
for $e=0,\ldots,7$ are
\[
 1,5,23,75,178,378,661,1065.
\]
All pushforwards include their residue-degree multiplicities.

The exhaustive source census at $5$ gives the following noncuspidal
closed-point counts in degrees $1$ through $10$:
\[
 0,0,16,64,656,2712,10784,48992,217520,977528.
\]
The corresponding diamond-orbit counts in degrees $3$ through $10$
are $2,8,82,342,1348,6172,27190,122318$.
These counts also agree with the independently computed Frobenius
polynomial. Since the least noncuspidal degree is $3$, the only mixed
support types are pairs of total degree at most $10$, and the triples
$(3,3,3)$ and $(3,3,4)$. Repeated supports are included.

For a single support of degree $9$ or $10$, four quadratic characters
annihilating $H$ give a preliminary form of (48.1). They are evaluated
using rational functions with even divisors; the functions, divisor-class
characters and base values are retained in
\path{inputs/four_characters.m}. A pole or zero of an evaluation is
retained for the full class test. Of $149508$ source diamond orbits,
$12604$ pass this preliminary test, including one pole fallback.
Their images give $5619$ queries after the rational automorphisms of $Y$.
The full class test leaves $107$ source cases: $57$ of degree $10$ and
$50$ of degree $9$ with a cusp complement of degree $1$.

For the lower and mixed supports, we retain every actual closed point
needed in a second or third summand; only a simultaneous source diamond
is used to normalize the whole divisor. The $16$ target automorphisms
are used to cache individual class queries. Their actions are undone
before summing mixed classes. The $3550$ target queries test all
$119986$ listed support patterns. Condition (48.1) leaves $243$ cases,
which give $201$ simultaneous source-diamond orbits.
Every class query is certified by an exact identity
$E-R-\operatorname{div}(f)$ in both maximal orders, or by an exact
factor-base hit. The independent replay also checks the actual point
coordinates and all norm weights.

It remains to exclude the resulting $308$ cases. For the $107$ cases
with no cusp complement ($57$ high and $50$ lower cases), Lemma~\ref{lem:decic48NC13} makes
$A_7N$ principal. For the other $201$ cases ($50$ high and $151$ lower cases), the same lemma makes
$A_{13}N$ principal, including every possible cuspidal complement.
We construct these divisors on the full genus-$37$ source, using all
cyclic $7$- or $13$-kernels with their Frobenius-orbit lengths.
For a source support of degree $d$ and a kernel orbit of length $e$,
a target prime of degree $f$ occurs with norm weight $de/f$.
The marked image retains exact order $48$.

In every one of the $308$ cases the resulting degree-zero divisor
has Riemann--Roch dimension $0$. All sums, including repeated supports,
are independently rebuilt in a fresh source function field. Separate
Sage computations enumerate all geometric cyclic kernels, partition
by Frobenius, evaluate normalized isogenies and independent pairwise
V\'elu formulas, and bind the marked source, diamond and target primes
to those used in the native Riemann--Roch calculations.
Thus none of the required principal divisors exists, a contradiction.
\end{proof}
